(1) 均匀面电荷密度 $\sigma=Q/(4\pi a^2)$，球壳表面线速度为 $\mathbf v=\boldsymbol\Omega\times\mathbf r=\Omega a\sin\theta\,\hat{\boldsymbol\varphi}$。因而面电流 $\mathbf K=\sigma\mathbf v=\sigma\Omega a\sin\theta\,\hat{\boldsymbol\varphi}$；$Q<0$ 时其方向反转。
(2) 偶极矩 $\boldsymbol{\mathfrak m}=\frac12\int\mathbf r\times\mathbf K\,dA$，取 $z$ 分量并使用 $\int\sin^2\theta\,dA=8\pi a^2/3$，得 $\mathfrak m=Q\Omega a^2/3$。面电流与均匀磁化球 $\mathbf M=\sigma\Omega a\,\hat{\mathbf z}$ 的束缚电流相同，故球内 $\mathbf B_{\rm in}=(2\mu_0/3)\mathbf M=\mu_0\boldsymbol{\mathfrak m}/(2\pi a^3)$；球外为偶极场 $\mathbf B_{\rm out}=(\mu_0/(4\pi r^3))[3(\boldsymbol{\mathfrak m}\cdot\hat{\mathbf r})\hat{\mathbf r}-\boldsymbol{\mathfrak m}]$。
(3) 球内场均匀，积分 $U_{\rm in}=\int_{r<a}B_{\rm in}^2/(2\mu_0)dV=\mu_0\mathfrak m^2/(6\pi a^3)$；球外偶极场平方的角积分给 $U_{\rm out}=\int_a^\infty\mu_0\mathfrak m^2(1+3\cos^2\theta)/(32\pi^2r^6)\,dV=\mu_0\mathfrak m^2/(12\pi a^3)$。总能量 $U_B=\mu_0\mathfrak m^2/(4\pi a^3)$。
(4) 磁偶极矩与电流均正比于 $Q\Omega$，磁能因而正比于 $Q^2\Omega^2$，在 $Q\to0$ 或 $\Omega\to0$ 时趋零。球外 $r\gg a$ 场严格按偶极 $r^{-3}$ 衰减；静电球壳仍有独立的库仑电场能 $U_E=Q^2/(8\pi\varepsilon_0 a)$，二者不能混同。
